% TEMPLATE-MILC-v3.tex - plantilla maestra UNICA de las guias MILC v3.
%
% La usan las cuatro familias de guias, cada una con su driver:
%   · Grados 6-11          -> scripts/build-guias-g11.py
%   · Bebras               -> scripts/build-guias-bebras.py
%   · Semillero            -> scripts/build-guias-semillero.py
%   · Territorio interior  -> scripts/build-guias-territorio-interior.py
%
% Antes habia tres copias casi identicas (~400 lineas iguales, 6 distintas):
% cada mejora habia que aplicarla tres veces, y en la practica no se hacia
% (el motor de recursos vivia solo en dos de ellas). Lo que varia entre
% familias esta parametrizado en 6 placeholders que cada driver llena
% (se nombran aqui SIN su sintaxis real: el builder reemplaza por texto plano,
% tambien dentro de los comentarios, y un valor multilinea rompe el preambulo):
%
%   HEADER_IZQ        encabezado izquierdo de pagina
%   HEADER_DER        encabezado derecho de pagina
%   PORTADA_ETIQUETA  rotulo sobre el numero grande (GRADO/MOMENTO/MODULO)
%   PORTADA_SERIE     linea de serie de la portada
%   PORTADA_ID        identificador de la guia en la portada
%   PORTADA_CIRCULO   el circulito de grado (°), o vacio si no aplica
%   PDF_TITULO        titulo en texto plano (sin \textbf ni comillas TeX) para
%                     los metadatos del PDF (/Title); lo produce lib_xelatex.texto_pdf
%
% Composicion (auditoria 2026-09-03): las bandas de estacion piden espacio
% (\Needspace*) y prohiben el salto justo despues (\nopagebreak), asi no quedan
% huerfanas al pie; las cajas largas (softbox, drawbox, infoband, puentebox) son
% `breakable`, asi una caja de media pagina ya no arrastra todo a la siguiente
% dejando la anterior al 40 %. Todo texto sobre mostaza o verde va en milcNegro
% (blanco sobre #E5B400 da 1,93:1 y sobre #58A923 2,95:1; ninguno pasa WCAG AA).
% El resto de claves las documenta content/guias/_SCHEMA.md. El script valida
% que no queden placeholders sin reemplazar antes de escribir el archivo final.

% Etiquetado PDF (latex-lab): /Lang, estructura y texto alternativo de las
% figuras, para lectores de pantalla. Debe ir ANTES de \documentclass.
\DocumentMetadata{lang=es-CO,tagging=on}
\documentclass[11pt,letterpaper]{article}
% headheight=24pt: la cabecera derecha lleva dos lineas (identidad de la guia
% y estacion actual). headsep se acorta en la misma medida para que el bloque
% de texto quede exactamente donde estaba con headheight=12pt.
\usepackage[margin=2.0cm,headheight=24pt,headsep=13pt]{geometry}
\usepackage[table]{xcolor}
\usepackage{fontspec}
\usepackage[spanish]{babel}
\usepackage{tikz}
% Librerías TikZ para los diagramas de las guías (bloque `recursos:`):
%  · babel  → babel en español vuelve activos caracteres como «>», y TikZ los usa
%             en sus opciones (>=stealth, ->). Sin esta librería, cualquier
%             diagrama con flechas revienta con «\language@active@arg> has an
%             extra }». Debe ir ANTES de que se dibuje nada.
%  · positioning        → sintaxis «below left=2pt and 3pt».
%  · decorations.pathreplacing → llaves ({brace}) para señalar rangos.
\usetikzlibrary{babel,positioning,decorations.pathreplacing}
\usepackage{graphicx}
% Circuitos: el trabajo de electrónica se hace hoy con micro:bit y sus sensores
% integrados (MakeCode, sin cableado), así que no se necesitan esquemáticos.
% Si algún día entran montajes con protoboard, basta descomentar esta línea:
% los diagramas .tex/.tikz declarados en `recursos:` ya se incrustan solos.
% \usepackage{circuitikz}
\usepackage[most]{tcolorbox}
\usepackage{tabularx}
\usepackage{array}
\usepackage{fancyhdr}
\usepackage{titlesec}
\usepackage[document]{ragged2e}  % bandera a la derecha en todo el cuerpo
\usepackage{enumitem}
\usepackage{microtype}
\usepackage{etoolbox}
\usepackage{needspace}
% hyperref al final del preambulo: metadatos del PDF (titulo, autor, idioma,
% familia), marcadores por estacion (\pdfbookmark) y enlaces sin recuadro.
\usepackage[hidelinks,
  pdftitle={Anestesias: lo que se usa para no sentir, y lo que se lleva a cambio},
  pdfauthor={PhD. Álvaro Cárdenas Orozco},
  pdfsubject={Territorio interior · Educación socioemocional · Grado 10° · Momento 6 · MILC v3},
  pdflang={es-CO},
  bookmarksopen=true,bookmarksnumbered=false]{hyperref}

% Helvetica Neue no trae la flecha U+2192, asi que cada «→» escrito literalmente
% en el YAML se compilaba a NADA, en silencio (462 flechas en 61 guias: rutas del
% tipo «paleta → tipografia → lockup» salian rotas). Mapeamos el caracter a su
% version matematica, que si tiene glifo. Asi el autor puede seguir escribiendo
% «→» comodamente en el YAML.
\usepackage{newunicodechar}
\newunicodechar{→}{\ensuremath{\rightarrow}}
% Mismo problema con los simbolos de marca y vinetas: el «✓» de «una palomita ✓»
% salia como cuadrito vacio en el PDF de todas las guias que lo usaban.
\usepackage{pifont}
\newunicodechar{✓}{\ding{51}}
\newunicodechar{✗}{\ding{55}}
\newunicodechar{✔}{\ding{52}}
\newunicodechar{•}{\textbullet}
\newunicodechar{≥}{\ensuremath{\geq}}
\newunicodechar{≤}{\ensuremath{\leq}}

\setmainfont{Helvetica Neue}
\setsansfont{Helvetica Neue}
\setlength{\parindent}{0pt}
\setlength{\parskip}{6pt}
% Sin particion de palabras: en 6.º-7.º una palabra cortada por guion al final
% de linea («Comuni-cacion») frena la lectura mucho mas que una linea corta.
\hyphenpenalty=10000
\exhyphenpenalty=10000
\pretolerance=2000
\tolerance=3000
\setlength{\emergencystretch}{3em}
\linespread{1.10}
\renewcommand{\arraystretch}{1.25}
\setlist{leftmargin=5mm,itemsep=1mm,topsep=1mm}
\newcolumntype{Y}{>{\RaggedRight\arraybackslash}X}

\definecolor{milcMagenta}{HTML}{E6007E}
\definecolor{milcVino}{HTML}{5A0038}
\definecolor{milcTurquesa}{HTML}{0093A5}
\definecolor{milcVerde}{HTML}{58A923}
\definecolor{milcMostaza}{HTML}{E5B400}
\definecolor{milcOcre}{HTML}{B86600}
\definecolor{milcNegro}{HTML}{241816}
\definecolor{milcGris}{HTML}{F7F7F4}
\definecolor{milcLinea}{HTML}{E8E2DD}
\definecolor{milcGradePrimary}{HTML}{0D9488}
\definecolor{milcGradeDark}{HTML}{115E59}
\definecolor{milcGradeSoft}{HTML}{CCFBF1}
\definecolor{milcGradeAccent}{HTML}{CFE7FF}
\definecolor{milcGradeText}{HTML}{FFFFFF}
\color{milcNegro}

\pagestyle{fancy}
\fancyhf{}
% Cabecera de dos lineas: identidad de la guia arriba y, a la derecha y en
% gris, la estacion en curso (\markright desde \milcestacion). Antes de la
% primera estacion la segunda linea va vacia; el \strut mantiene la altura
% para que la linea de arriba no baile entre paginas.
\lhead{\small Territorio interior · Educación socioemocional\\[-1pt]{\footnotesize\strut}}
\rhead{\small Grado 10° · Momento 6 · MILC v3\\[-1pt]{\footnotesize\strut\color{milcNegro!65}\rightmark}}
\cfoot{\small\thepage}
\renewcommand{\headrulewidth}{0pt}

% Sobre mostaza (#E5B400) y verde (#58A923) el texto va en milcNegro; sobre el
% resto de la paleta, en blanco. Se usa dentro de fonttitle/fontupper, donde un
% \color posterior manda sobre coltitle.
\newcommand{\milctintasobre}[1]{%
  \ifboolexpr{test{\ifstrequal{#1}{milcMostaza}} or test{\ifstrequal{#1}{milcVerde}}}%
    {\color{milcNegro}}{\color{white}}}

\newtcolorbox{titlebox}[1]{
  enhanced,colback=#1,colframe=#1,arc=7mm,boxrule=0pt,
  left=7mm,right=7mm,top=3mm,bottom=3mm,
  before skip=8pt,after skip=8pt,
  fontupper=\sffamily\bfseries\Large\color{white}
}

\newtcolorbox{softbox}[3]{
  enhanced,breakable,pad at break=3mm,
  colback=#2,colframe=#1,borderline west={4pt}{0pt}{#1},
  arc=2mm,boxrule=.4pt,left=4mm,right=4mm,top=3mm,bottom=3mm,
  before skip=6pt,after skip=8pt,title={#3},coltitle=white,
  colbacktitle=#1,fonttitle=\sffamily\bfseries\milctintasobre{#1},fontupper=\RaggedRight,
  attach boxed title to top left={xshift=3mm,yshift=-2mm},
  boxed title style={arc=2mm, boxrule=0pt}
}

\newtcolorbox{drawbox}[2]{
  enhanced,breakable,pad at break=4mm,
  colback=white,colframe=#1,arc=4mm,boxrule=1pt,
  left=5mm,right=5mm,top=4mm,bottom=4mm,before skip=8pt,after skip=8pt,
  title={#2},coltitle=white,colbacktitle=#1,fonttitle=\sffamily\bfseries\milctintasobre{#1},
  fontupper=\RaggedRight,
  attach boxed title to top left={xshift=4mm,yshift=-2mm},
  boxed title style={arc=3mm, boxrule=0pt}
}

\newtcolorbox{infoband}[1]{
  enhanced,breakable,pad at break=2mm,colback=#1!8,colframe=#1!50,arc=2mm,boxrule=.5pt,
  left=4mm,right=4mm,top=2mm,bottom=2mm,before skip=4pt,after skip=4pt,
  fontupper=\small\RaggedRight
}

% Puente narrativo entre secciones (contrato editorial Conéctate).
% Da continuidad: "Acabas de X. Ahora vas a Y porque Z."
% Visualmente: banda izquierda en color del grado, texto en itálica pequeña.
\newtcolorbox{puentebox}{
  enhanced,breakable,pad at break=2mm,colback=milcGradeSoft,colframe=milcGradeDark,
  borderline west={3pt}{0pt}{milcGradeDark},
  arc=0mm,boxrule=0pt,left=5mm,right=4mm,top=2.5mm,bottom=2.5mm,
  before skip=4pt,after skip=8pt,
  fontupper=\itshape\small\color{milcNegro}\RaggedRight
}
% La flecha va en modo matematico a proposito: Helvetica Neue NO tiene el glifo
% U+2192, asi que \textrightarrow se compilaba a NADA («Missing character: There
% is no → in font Helvetica Neue») y el puente salia sin flecha. En modo
% matematico la flecha viene de la fuente matematica y si se dibuja.
\newcommand{\puente}[1]{%
  \begin{puentebox}%
    \textcolor{milcGradeDark}{$\rightarrow$}\hspace{2mm}#1%
  \end{puentebox}%
}

\newcommand{\linead}[1]{\noindent\color{milcLinea}\rule{#1}{0.5pt}\color{milcNegro}}
\newcommand{\respuesta}[1]{\par\vspace{2mm}\foreach \n in {1,...,#1}{\noindent\color{milcLinea}\rule{\linewidth}{0.5pt}\par\vspace{5mm}}\color{milcNegro}}
\newcommand{\checkbox}{\textcolor{milcGradeDark}{\fbox{\phantom{x}}}\hspace{5pt}}

% ─── Listas aireadas para las guías socioemocionales ────────────────────
% Componer las verificaciones como «\checkbox texto\\» deja el texto que salta
% de línea debajo del cuadrito y apelmaza el bloque. Estos entornos alinean la
% continuación y airean los ítems. Los usa Territorio Interior; cualquier
% familia puede hacerlo.
\newlist{checklista}{itemize}{1}
\setlist[checklista]{
  label=\textcolor{milcGradeDark}{\fbox{\phantom{x}}},
  leftmargin=9mm, labelsep=3.2mm, itemsep=2.4mm,
  topsep=2.5mm, parsep=0pt, partopsep=0pt, before=\RaggedRight
}
\newlist{pasoslista}{enumerate}{1}
\setlist[pasoslista]{
  label=\textcolor{milcGradeDark}{\sffamily\bfseries\arabic*.},
  leftmargin=8mm, labelsep=3mm, itemsep=2.8mm,
  topsep=1mm, parsep=0pt, partopsep=0pt, before=\RaggedRight
}

\newcommand{\phasecard}[4]{
\begin{tcolorbox}[enhanced,colback=#3,colframe=#2,arc=3mm,boxrule=.5pt,height=3.05cm,valign=center,halign=center,left=1.5mm,right=1.5mm,top=2mm,bottom=2mm]
{\sffamily\bfseries\fontsize{22}{24}\selectfont\textcolor{#2}{#1}}\par
{\sffamily\bfseries\small #4}
\end{tcolorbox}
}

% ── Piezas del rediseno 2026 (legibilidad 6.º-11.º) ───────────────────
% Banda de estacion: reemplaza el titlebox macizo. Titulo a la izquierda,
% dato practico (tiempo/modalidad) a la derecha, en una sola linea.
% Nunca huerfana: pide 9 lineas libres (o salta de pagina rellenando con
% \vfill) y prohibe el salto justo despues de la banda. Nueve y no seis:
% la caja partible que sigue necesita ~80pt (titulo adosado, relleno y dos
% lineas) para arrancar en la misma pagina; con menos, tcolorbox la manda
% entera a la siguiente y la banda se queda sola. El penalty 10000 va
% en `after`, ANTES del espacio posterior: un `after skip` (aunque sea 0pt)
% es un \vskip tras la caja, y ese glue es un punto de corte legal que un
% \nopagebreak escrito despues de \end{tcolorbox} ya no cierra. Cada banda
% deja marcador PDF y actualiza la cabecera derecha.
\newcounter{milcestacion}
\newcommand{\milcestacion}[2]{%
\Needspace*{9\baselineskip}%
\stepcounter{milcestacion}%
\pdfbookmark[1]{#2}{milcestacion\themilcestacion}%
\markright{#2}%
\begin{tcolorbox}[enhanced,colback=#1,colframe=#1,arc=2mm,boxrule=0pt,
  left=5mm,right=5mm,top=2.6mm,bottom=2.6mm,before skip=11pt,
  after={\par\nopagebreak\vskip7pt}]
{\sffamily\bfseries\fontsize{15}{18}\selectfont\milctintasobre{#1}#2}
\end{tcolorbox}}

% Tarjeta de paso numerada: sustituye las tablas «Pilar / Aplicacion», que
% eran formato de consulta docente, no de estudiante. El numero va en un
% circulo sobre el borde para que la secuencia se lea de un vistazo.
\newcommand{\milcpaso}[3]{%
\Needspace{4\baselineskip}%
\begin{tcolorbox}[enhanced,colback=white,colframe=milcLinea,arc=1.5mm,boxrule=.9pt,
  left=14mm,right=4mm,top=3mm,bottom=3mm,before skip=4pt,after skip=4pt,
  fontupper=\RaggedRight,
  overlay={\node[anchor=north west,circle,fill=#2,text=white,inner sep=0pt,
    minimum size=7.4mm,font=\sffamily\bfseries\small]
    at ([xshift=3.6mm,yshift=-3.2mm]frame.north west) {#1};}]
#3
\end{tcolorbox}}

% Casilla marcable de tamano util con lapiz (la anterior era \fbox{\phantom{x}},
% de alto variable segun el texto que la seguia).
\newcommand{\milcmarca}{\framebox[3.8mm]{\rule{0pt}{3.4mm}}}

% Fila marcable de lista suelta (checks de la estacion 1).
\newcommand{\milccheckrow}[1]{%
\par\vspace{1.8mm}\noindent
\makebox[6.4mm][l]{\raisebox{-.55mm}{\textcolor{milcNegro}{\milcmarca}}}%
\parbox[t]{\dimexpr\linewidth-6.4mm\relax}{\RaggedRight #1}\par}

% Celda marcable dentro de la rubrica (tabla de autoevaluacion). Misma
% sangria francesa, pero pensada para vivir dentro de una columna X.
\newcommand{\milcopcion}[1]{%
\makebox[5.2mm][l]{\raisebox{-.55mm}{\textcolor{milcNegro}{\milcmarca}}}%
\parbox[t]{\dimexpr\linewidth-5.2mm\relax}{\RaggedRight #1}}

% ── Recursos visuales de la guía ──────────────────────────────────────
% Declarados en el bloque `recursos:` del YAML y emitidos por el builder.
% \guiaFigura   → imagen rasterizada o PDF (foto, carta celeste, montaje).
% \guiaDiagrama → fuente TikZ/circuitikz (.tex) incrustada como vector:
%                 es la vía para circuitos y esquemáticos editables.
% [#1] = texto alternativo (alt) · #2 = ruta relativa al .tex · #3 = pie
% El alt llega al PDF etiquetado (/Alt) y lo lee el lector de pantalla.
\newcommand{\guiaFigura}[3][]{%
\begin{center}
\includegraphics[width=0.88\linewidth,height=0.38\textheight,keepaspectratio,alt={#1}]{#2}\\[1.5mm]
{\footnotesize\itshape\textcolor{milcNegro!75}{#3}}
\end{center}
}

\newcommand{\guiaDiagrama}[2]{%
\begin{center}
\input{#1}\\[1.5mm]
{\footnotesize\itshape\textcolor{milcNegro!75}{#2}}
\end{center}
}

\begin{document}
\thispagestyle{empty}

% ── Portada ───────────────────────────────────────────────────────────
% Antes: un rectangulo de color a pagina completa. Se veia bien en pantalla,
% pero en la fotocopiadora del colegio se comia el toner y salia gris sucio.
% Ahora la banda ocupa el tercio superior y el resto va en blanco.
\begin{tikzpicture}[remember picture,overlay]
  \path[fill=milcGradePrimary]
    (current page.north west) rectangle ([yshift=-10.6cm]current page.north east);

  \node[anchor=north west,text=milcGradeText] at ([xshift=2.0cm,yshift=-1.55cm]current page.north west)
    {\sffamily\bfseries\fontsize{12}{14}\selectfont MOMENTO};
  \node[anchor=north west,text=milcGradeText,opacity=.85] at ([xshift=2.0cm,yshift=-2.35cm]current page.north west)
    {\sffamily\bfseries\fontsize{8.5}{10}\selectfont TERRITORIO INTERIOR · SOCIOEMOCIONAL};
  \node[anchor=north east,text=milcGradeText,opacity=.85,align=right] at ([xshift=-2.0cm,yshift=-1.55cm]current page.north east)
    {\sffamily\bfseries\fontsize{9}{12}\selectfont I.E. Sor María Juliana\\Cartago, Valle del Cauca};

  \node[anchor=west,text=milcGradeText] (gradonum) at ([xshift=1.85cm,yshift=-7.1cm]current page.north west)
    {\sffamily\bfseries\fontsize{112}{112}\selectfont 6};
  % El circulito de grado (°) solo aplica a las guias de grado («11°»); en
  % Bebras y Semillero el numero grande es un momento/modulo, no un grado.
  % Se posiciona relativo al nodo `gradonum`, no en coordenadas absolutas.
  

  \node[anchor=west,text=milcGradeText,text width=11.4cm,align=left] at ([xshift=1.5cm,yshift=0.5cm]gradonum.east)
    {
     {\sffamily\bfseries\fontsize{9}{11}\selectfont\MakeUppercase{Territorio interior · Socioemocional}}\\[3mm]
     {\sffamily\bfseries\fontsize{21}{25}\selectfont\setlength{\baselineskip}{27pt}Anestesias\\[4pt]lo que se usa\\[4pt]para no sentir}\\[3.5mm]
     {\sffamily\bfseries\fontsize{15}{18}\selectfont Grado 10° · Momento 6}
    };
\end{tikzpicture}

\vspace*{9.4cm}
\vfill

{\sffamily\bfseries\fontsize{8}{10}\selectfont\textcolor{milcGradeDark}{LO QUE TE LLEVAS HOY}}

\begin{tcolorbox}[enhanced,colback=white,colframe=milcNegro,arc=2mm,boxrule=1.6pt,
  left=5mm,right=5mm,top=4mm,bottom=4mm,before skip=3pt,after skip=12pt,
  fontupper=\RaggedRight\fontsize{11}{15.5}\selectfont]
Tu mapa de anestesias: qué usas cuando algo se pone difícil y con qué frecuencia, la emoción que aparece justo antes de cada una, lo que has dejado de hacer desde que las usas, y un plan de sustitución con una alternativa concreta para las tres situaciones más frecuentes.
\end{tcolorbox}

\begin{tcolorbox}[enhanced,colback=milcGris,colframe=milcGris,arc=2mm,boxrule=0pt,
  left=5mm,right=5mm,top=3.5mm,bottom=3.5mm,after skip=12pt]
{\sffamily\bfseries\fontsize{8}{10}\selectfont\textcolor{milcNegro!55}{ESTA GUÍA ES DE}}\\[3mm]
\begin{tabularx}{\linewidth}{X p{3.2cm} p{3.2cm}}
\linead{\linewidth} & \linead{3.0cm} & \linead{3.0cm}\\[-1mm]
{\fontsize{8.5}{10}\selectfont\textcolor{milcNegro!55}{Nombre completo}} &
{\fontsize{8.5}{10}\selectfont\textcolor{milcNegro!55}{Grupo}} &
{\fontsize{8.5}{10}\selectfont\textcolor{milcNegro!55}{Fecha}}\\
\end{tabularx}
\end{tcolorbox}

\vfill

\noindent\textcolor{milcLinea}{\rule{\linewidth}{1pt}}\\[2mm]
{\sffamily\bfseries\fontsize{9.5}{12}\selectfont Autor: Dr. Álvaro Cárdenas Orozco}\\[1mm]
{\sffamily\fontsize{8.5}{11}\selectfont\textcolor{milcNegro!55}{Metodología MILC · Apertura ancestral · Evitación experiencial · Triángulo Dussel-Estoicismo-Floridi}}

\newpage

\section*{Apertura: Anestesias: lo que se usa para no sentir, y lo que se lleva a cambio}

\begin{softbox}{milcGradeDark}{milcGradeSoft}{Saber ancestral}
En los años noventa llegó a \textbf{Tierradentro}, territorio nasa en el Cauca, un cultivo nuevo: la \textbf{amapola} ---la planta de la que salen el opio y la heroína---. \emph{Prometía plata rápida, y algunos cultivadores la sembraron.} \textbf{Lo que pasó después está documentado, y es más sutil de lo que uno esperaría.} \emph{No hubo un colapso repentino:} \textbf{el cultivo alteró las dinámicas de la comunidad y las prácticas culturales, y lo primero que cambió fue lo más cotidiano ---\emph{la forma de vestir y la forma de comer}---.} \textbf{Los estudios lo describen sin rodeos:} \emph{la amapola llegó a constituir \textbf{una amenaza real para la continuidad cultural} de los nasa de Tierradentro}. \emph{Fíjate en el orden de los hechos, porque es exactamente el de esta guía:} \textbf{lo que entra prometiendo alivio no anuncia lo que se va a llevar, y lo primero que se lleva son cosas pequeñas que nadie estaba mirando.} \textbf{Y hay un final que importa:} \emph{cuando ocurrió la avalancha del \textbf{6 de junio de 1994}, los médicos tradicionales la leyeron como un reclamo de la madre tierra ---agredida por la deforestación y por los cultivos---, y \textbf{la comunidad dejó de sembrar amapola}.} \emph{No los detuvo una prohibición: los detuvo darse cuenta de lo que estaban perdiendo.}
\end{softbox}

\begin{softbox}{milcTurquesa}{milcGris}{Saber contemporáneo}
¿Qué tienen en común el trago, las apuestas, la pantalla a las tres de la mañana y comer sin hambre? \textbf{Que hacen lo mismo.} Un grupo de investigadores propuso mirar buena parte del sufrimiento psicológico desde una sola idea: la \textbf{evitación experiencial}, es decir, \emph{los intentos ---poco saludables--- de \textbf{escapar y evitar} las emociones propias}. \textbf{Revisando una cantidad grande de investigación, sostuvieron que muchas formas de malestar pueden entenderse así:} \emph{no como problemas distintos con sustancias distintas, sino como \textbf{la misma maniobra} aplicada con herramientas distintas} (Hayes, Wilson, Gifford, Follette y Strosahl, 1996). \textbf{Eso reordena por completo la conversación.} \emph{La pregunta útil deja de ser «¿qué está usando esta persona?» y pasa a ser «\textbf{¿qué está tratando de no sentir?}».} \textbf{Y explica algo que se ve todo el tiempo:} \emph{cuando alguien deja una anestesia sin haber resuelto lo que evitaba, casi siempre aparece otra} ---deja de fumar y empieza a apostar, deja de apostar y no suelta el teléfono---. \textbf{Porque el problema nunca fue la herramienta.}
\end{softbox}

\begin{softbox}{milcMostaza}{milcGris}{Pregunta-puente}
\textit{En Tierradentro, lo primero que cambió la amapola \textbf{fue la forma de vestir y de comer} ---nada que pareciera grave---. \textbf{Cuando algo te angustia, ¿qué es lo primero que haces} \ldots \textbf{y qué has dejado de hacer desde que lo haces}?}
\end{softbox}

\begin{softbox}{milcVerde}{milcGris}{Saber y saber hacer}
Al terminar podrás: (1) \textbf{reconocer} que sustancias, apuestas y pantallas cumplen la misma función ---evitar una emoción--- y que por eso se reemplazan entre sí; (2) \textbf{identificar} qué emoción aparece justo \emph{antes} de tu propia anestesia; (3) \textbf{ver} el costo donde de verdad se paga: en lo que has dejado de hacer; (4) \textbf{planear} una sustitución concreta para tus tres situaciones más frecuentes, y saber cuándo esto ya no se maneja solo.
\end{softbox}



\small
\begin{tabularx}{\textwidth}{>{\bfseries}p{3.0cm}Y}
\rowcolor{milcGradePrimary!12}Autor & Dr. Álvaro Cárdenas Orozco\\
DBA & Comprende los fenómenos que afectan a su territorio, regula sus emociones ante ellos y acompaña a otros con empatía (transversal 6°--11°, articulado con la política «Escuela, territorio de vida» del MEN, Resolución 006519 de 2025, y la Ley 1620 de 2013)\\
Referentes & Primera ayuda psicológica (OMS, 2011/2012) · Directrices IASC (2007) · USGS y Servicio Geológico Colombiano · Pueblo nasa y la minga (CRIC) · Dussel · Estoicismo · Floridi\\
Producto final & Tu mapa de anestesias: qué usas cuando algo se pone difícil y con qué frecuencia, la emoción que aparece justo antes de cada una, lo que has dejado de hacer desde que las usas, y un plan de sustitución con una alternativa concreta para las tres situaciones más frecuentes.\\
\end{tabularx}
\normalsize

\puente{En el momento 1 viste que \textbf{evitar alivia hoy y agrava mañana}. \textbf{Hoy le pones nombre a las herramientas} con las que se evita, y son más de las que uno cree.}

\milcestacion{milcGradeDark}{Ruta MILC: así vamos a aprender}
\begin{tabularx}{\textwidth}{>{\centering\arraybackslash}X>{\centering\arraybackslash}X>{\centering\arraybackslash}X>{\centering\arraybackslash}X}
\phasecard{1}{milcVerde}{milcGris}{Escucha\\[-1mm]\small Observo, escucho y pregunto.} &
\phasecard{2}{milcTurquesa}{milcGris}{Sistematización\\[-1mm]\small Veo qué estoy evitando.} &
\phasecard{3}{milcMagenta}{milcGris}{Praxis\\[-1mm]\small Construyo y aplico.} &
\phasecard{4}{milcVino}{milcGris}{Evaluación\\[-1mm]\small Reflexión liberadora.}
\end{tabularx}


\puente{Primero, sin juzgarte: qué haces cuando algo se pone difícil. \emph{Incluye lo que no parece grave, que es casi todo.}}

\milcestacion{milcVerde}{Estación 1 de 3 · Escucha}

\begin{softbox}{milcVerde}{milcGris}{Escena para escuchar}
\textbf{Actividad 1 · IDENTIFICA --- Qué hago cuando se pone difícil.} \textbf{Nadie va a leer esta página, y aquí no se juzga a nadie ---tampoco a ti---.} \textbf{Primera parte.} Escribe \textbf{qué haces cuando algo te angustia, te aburre o te duele}. \emph{Todo, incluido lo que no parece grave:} \textbf{el teléfono, las series, dormir de más, comer, no comer, videojuegos, apuestas, trago, cigarrillo, vapeador, encerrarte, pelear, no salir de la cama.} \textbf{Segunda parte.} Al frente de cada uno: \emph{¿cuántas veces esta semana? ¿cuánto tiempo se te va?} \textbf{Tercera parte, la difícil y la más útil.} \emph{¿Qué estabas sintiendo \textbf{justo antes}?} \textbf{Usa el vocabulario del grado 6, momento 2:} \emph{no «mal», sino «aburrido», «solo», «con rabia», «avergonzado», «asustado», «cansado».} \textbf{Y la última:} \emph{¿hay algo que hayas dejado de hacer} ---dormir bien, estudiar, ver a alguien, un deporte, salir--- \emph{desde que haces eso?}
\end{softbox}

{\sffamily\bfseries\fontsize{8}{10}\selectfont\textcolor{milcNegro!55}{MARCA CUANDO LO HAYAS HECHO}}
\milccheckrow{Listaste todo, incluido lo que no parece grave.}
\milccheckrow{Anotaste frecuencia y tiempo de cada uno.}
\milccheckrow{Escribiste \textbf{qué sentías justo antes}, con una palabra precisa.}

\begin{infoband}{milcVerde}


\textbf{En tu cuaderno (Actividad 1 · IDENTIFICA).} Título: «Actividad 1. Qué hago cuando se pone difícil». \textbf{Formato:} la lista con frecuencia y tiempo + la emoción de antes + lo que has dejado de hacer. \textbf{Extensión:} media página. \textbf{Sabes que terminaste cuando} tienes \textbf{una emoción escrita} al frente de cada una. Esta página es tuya: \textbf{nadie más tiene que leerla si tú no quieres}.
\end{infoband}

\puente{Ya lo tienes. Ahora, por qué todas hacen lo mismo, por qué funcionan ---que es el problema--- y dónde se paga el costo.}

\milcestacion{milcTurquesa}{Fase 2 · Sistematización: entender la anestesia}

\begin{softbox}{milcTurquesa}{milcGris}{Cuatro claves sobre no sentir}
Cuatro claves para entender algo que no se arregla con fuerza de voluntad.
\end{softbox}

\milcpaso{1}{milcTurquesa}{\textbf{Todas hacen lo mismo, y por eso se reemplazan entre sí.} \emph{La idea de la \textbf{evitación experiencial} es que mucho malestar se sostiene en intentos poco saludables de \textbf{escapar y evitar} las emociones propias} (Hayes y otros, 1996). \textbf{Bajo esa lente, el trago, la apuesta, la pantalla, la comida y encerrarse no son cinco problemas: son \emph{una misma maniobra} con cinco herramientas.} \emph{Y eso explica lo que se ve todo el tiempo:} \textbf{quien deja una sin resolver lo que evitaba, encuentra otra} ---deja de fumar y empieza a apostar, deja de apostar y no suelta el teléfono---. \textbf{Por eso la pregunta no es «¿qué usa?» sino «\emph{¿qué está tratando de no sentir?}».}}
\milcpaso{2}{milcTurquesa}{\textbf{Funciona, y ese es exactamente el problema.} \emph{Sería fácil si no sirviera ---nadie repite lo que no sirve---.} \textbf{La anestesia \emph{sí} quita el malestar, y lo quita rápido:} \emph{por eso el cerebro aprende, en cuestión de días, que ese es el camino corto.} \textbf{Y ahí empieza el círculo del momento 1:} \emph{alivio inmediato, y la comprobación de que se podía tolerar \textbf{nunca llega}}. \emph{Con el tiempo pasan dos cosas:} \textbf{hace falta más para el mismo alivio}, y \textbf{el malestar de base crece}, porque lo que lo causaba sigue ahí sin tocarse. \emph{Por eso «tengo fuerza de voluntad» no es una defensa:} \textbf{esto no se sostiene en la debilidad de nadie ---se sostiene en que funciona a corto plazo---.}}
\milcpaso{3}{milcTurquesa}{\textbf{Se paga con lo que se deja, y ahí es donde hay que mirar.} \emph{Casi nadie nota una anestesia por su cantidad ---todo el mundo cree usar menos de lo que usa---.} \textbf{Se nota por el hueco:} \emph{qué dejaste de hacer}. \textbf{Y en Tierradentro pasó igual, a escala de pueblo:} \emph{lo primero que cambió la amapola fue \textbf{la forma de vestir y de comer}}, \textbf{y solo después se vio que estaba en juego la continuidad cultural.} \emph{Nadie decidió cambiar eso ---se fue yendo---.} \textbf{Tradúcelo:} \emph{el costo no aparece en las horas de pantalla ---aparece en el deporte que dejaste, en la persona que no ves, en las noches que no duermes, en lo que ya no te da gusto---.}}
\milcpaso{4}{milcTurquesa}{\textbf{Llega prometiendo otra cosa.} \emph{La amapola no llegó anunciando que iba a amenazar la continuidad cultural de nadie:} \textbf{llegó prometiendo plata rápida en un territorio donde la plata faltaba.} \emph{Y así llegan todas:} \textbf{el vapeador llega como algo inofensivo y de sabores; la apuesta llega como una forma de ganarse algo; la pantalla llega como descanso.} \textbf{Ninguna se presenta diciendo lo que va a cobrar.} \emph{Y hay algo más, incómodo y verdadero:} \textbf{casi siempre llegan donde hace falta algo} ---plata, compañía, sueño, sentido---. \emph{Por eso quitar sin poner no funciona: si algo llenaba un hueco real, hay que resolver el hueco.}}

\begin{drawbox}{milcTurquesa}{Cómo se arma una sustitución que sí funcione}
\begin{pasoslista}
\item \textbf{La situación:} \emph{¿cuándo pasa?} ---la hora, el lugar, el día de la semana---.
\item \textbf{La emoción de antes:} \emph{una palabra precisa}, no «mal».
\item \textbf{Qué te da:} \emph{calma, compañía, distracción, dejar de pensar, sentirte capaz.}
\item \textbf{Con qué otra cosa se consigue \emph{eso mismo}:} \emph{y tiene que dar lo mismo, no algo distinto.}
\item \textbf{Que esté a la mano:} \emph{una alternativa que hay que preparar veinte minutos no se usa nunca.}
\end{pasoslista}
\end{drawbox}

\begin{softbox}{milcMostaza}{milcGris}{Errores comunes y su corrección}
\textbf{«Es que no tengo fuerza de voluntad»:} no se sostiene en debilidad, sino en que funciona rápido. \textbf{Quitar sin poner:} el hueco se llena solo, y casi siempre con otra cosa igual. \textbf{Contar solo la cantidad:} el costo se ve en lo que dejaste de hacer. \textbf{Juzgar al que usa:} garantiza que no cuente nada. \textbf{Creer que solo cuentan las sustancias:} la pantalla y las apuestas hacen lo mismo. \textbf{Elegir una alternativa que no da lo mismo:} si buscabas compañía, salir a correr solo no sirve. \textbf{Pensar que esto siempre se resuelve solo:} a veces no, y eso se cuenta.
\end{softbox}

\begin{infoband}{milcTurquesa}


\textbf{En tu cuaderno (Actividad 2 · EXPLICA).} Título: «Actividad 2. La misma maniobra». \textbf{Formato:} explica con tus palabras qué es la evitación experiencial y por qué quien deja una anestesia suele encontrar otra. \textbf{Extensión:} media página. \textbf{Sabes que terminaste cuando} puedes explicar por qué \textbf{el costo se mide en lo que se dejó de hacer}.
\end{infoband}

\puente{Ya sabes el mecanismo. Es hora de lo que sí funciona: no prohibirse nada, sino \textbf{poner otra cosa} en el lugar exacto.}

\milcestacion{milcMagenta}{Fase 3 · Praxis: sustituir, no prohibir}

\begin{softbox}{milcMagenta}{milcGris}{Cómo se cambia una anestesia por otra cosa}
Aquí no se prohíbe nada. Se busca qué está haciendo esa anestesia por ti, y qué otra cosa podría hacer lo mismo.
\end{softbox}

\milcpaso{1}{milcMagenta}{\textbf{El registro de la semana (7 días, 2 min al día).} Anota \textbf{cada vez} que uses algo de tu lista, con \textbf{hora y situación}. \emph{Sin juzgarte y sin cambiar nada todavía ---solo mirar---.} \textbf{Al final, cuenta.} \emph{Casi todo el mundo se lleva la misma sorpresa, y la sorpresa es el punto: son más veces de las que creía, y casi siempre a la misma hora.}}
\milcpaso{2}{milcMagenta}{\textbf{La emoción de antes (8 min).} Al frente de cada registro, la palabra precisa: \emph{aburrido, solo, con rabia, avergonzado, asustado, cansado, sin sentido}. \textbf{Busca el patrón:} \emph{casi siempre son una o dos emociones las que disparan casi todo}, y saber cuáles son cambia por completo dónde hay que trabajar.}
\milcpaso{3}{milcMagenta}{\textbf{Lo que dejé de hacer (7 min).} Escribe la lista del hueco: \emph{qué hacías antes y ya no}. \textbf{Deporte, dormir, una amistad, dibujar, leer, salir, estudiar sin interrupciones.} \emph{Esta lista es el costo real, y es la que convence ---cuando convence--- porque la escribiste tú.}}
\milcpaso{4}{milcMagenta}{\textbf{El plan de sustitución (12 min).} Toma \textbf{tus tres situaciones más frecuentes} y para cada una escribe: \emph{la hora, la emoción, \textbf{qué te da} esa anestesia y \textbf{qué otra cosa podría darte lo mismo}}. \textbf{Y dos condiciones que deciden si funciona:} \emph{que la alternativa dé \textbf{lo mismo}} ---si lo que buscabas era compañía, la alternativa tiene que traer gente--- \emph{y que esté \textbf{a la mano}} ---preparada de antemano, no algo que haya que organizar en el momento---. \textbf{Y una nota que no es de relleno:} \emph{si algo de tu lista lleva meses, aumentó de cantidad, o ya intentaste dejarlo y no pudiste, \textbf{eso no se maneja solo}}: \textbf{es la ruta del momento 3, y usarla no es un fracaso.}}

\begin{drawbox}{milcMagenta}{Antes de cerrar tu mapa}
\begin{checklista}
\item Registraste \textbf{siete días} con hora y situación, sin cambiar nada.
\item Cada registro tiene una \textbf{emoción precisa} al frente, no «mal».
\item Encontraste \textbf{cuál es la emoción que dispara casi todo}.
\item Escribiste la lista de \textbf{lo que dejaste de hacer}.
\item Tus tres alternativas dan \textbf{lo mismo} que la anestesia y están \textbf{a la mano}.
\item Sabes cuándo esto \textbf{deja de manejarse solo} y cuál es tu ruta.
\end{checklista}
\end{drawbox}

\begin{softbox}{milcMagenta}{milcGris}{Plantilla de trabajo}
\textbf{El registro.} \emph{«Hora \_\_\_\_ · situación \_\_\_\_ · antes sentía \_\_\_\_.»} \\[1mm]
\textbf{La sustitución.} \emph{«Cuando \_\_\_\_, lo que busco es \_\_\_\_. En vez de \_\_\_\_ voy a \_\_\_\_, que ya tengo listo.»} \\[1mm]
\textbf{El límite.} \emph{«Si lleva meses, si aumentó o si ya intenté y no pude → hablo con \_\_\_\_.»}
\end{softbox}

\begin{infoband}{milcMagenta}


\textbf{En tu cuaderno (Actividad 3 · CREA).} Título: «Actividad 3. Mi mapa de anestesias». \textbf{Formato:} el registro de siete días con el conteo + la emoción de cada uno y el patrón + la lista del hueco + el plan de sustitución para tres situaciones. \textbf{Extensión:} una página. \textbf{Sabes que terminaste cuando} tus tres alternativas \textbf{están preparadas} ---no solo escritas---.
\end{infoband}

\puente{El producto es un plan de sustitución. \emph{Quitar sin poner deja un hueco, y el hueco se llena solo.}}

\milcestacion{milcMostaza}{Producto de la sesión}
\textbf{Mapa de anestesias: registro de siete días, la emoción que dispara y un plan de sustitución para las tres situaciones más frecuentes}.
\begin{softbox}{milcMostaza}{milcGris}{Criterios de calidad}
(1) El registro cubre siete días con hora y situación, sin intentar cambiar la conducta. (2) Cada episodio tiene identificada la emoción previa con una palabra precisa. (3) Se identifica el patrón: qué una o dos emociones disparan la mayoría de los episodios. (4) El costo está formulado como lo que se dejó de hacer, no como cantidad consumida. (5) Las alternativas cumplen la misma función que la anestesia y están disponibles de antemano.
\end{softbox}

\puente{Antes de cerrar, revisa lo que casi nadie escribe: si no identificaste la emoción de antes, el plan va a fallar.}

\milcestacion{milcVino}{Evaluación liberadora · cinco dimensiones}

\begin{softbox}{milcVino}{milcGris}{Sentido de la evaluación}
Evaluar no es solo revisar una nota. Es reconocer qué aprendí, qué cambió en mí, qué emociones manejé, cómo me relaciono con la ciudad y la comunidad, y qué le dejaré a quien viene detrás.
\end{softbox}

\textbf{Mi reflexión liberadora en cinco dimensiones.} Completa cada frase iniciada:

\small
\begin{tabularx}{\textwidth}{>{\bfseries\RaggedRight\arraybackslash}p{3.4cm}Y>{\RaggedRight\arraybackslash}p{4.6cm}}
\rowcolor{milcVino}\color{white}Dimensión & \color{white}\bfseries Mi respuesta (frame starter) & \color{white}\bfseries Algo que descubrí\\
1. Personal & Lo que más cambió en mí fue \linead{6.5cm} & \linead{4.2cm}\\
2. Emocional & La emoción más fuerte fue \linead{2.5cm} y la regulé \linead{3cm} & \linead{4.2cm}\\
3. Ciudadana & En democracia esto importa porque \linead{6cm} & \linead{4.2cm}\\
4. Local (Cartago) & En mi barrio sería útil para \linead{6.5cm} & \linead{4.2cm}\\
5. Intergeneracional & A mi abuela/abuelo le diría \linead{6.5cm} & \linead{4.2cm}\\
\end{tabularx}
\normalsize

\begin{infoband}{milcVino}
\textbf{Para la próxima sustentación.} Vuelve a esta tabla en seis meses. Verás que las respuestas cambian — no porque lo aprendido cambie, sino porque tú cambias. La evaluación liberadora es una conversación contigo a través del tiempo.
\end{infoband}


\puente{Cerramos con tres voces: el que está más allá de la totalidad, la demora como remedio, y el entorno que aprende de lo que miras.}

\milcestacion{milcGradeDark}{Triángulo de pensamiento · cierre filosófico}

\begin{softbox}{milcVino}{milcGris}{Dussel · lente del nosotros}
\textit{«Analéctico quiere indicar el hecho real humano por el que todo hombre, todo grupo o pueblo se sitúa siempre más allá (aná-) del horizonte de la totalidad.»}

{\footnotesize\itshape\color{milcNegro!70}--- Enrique Dussel · Filosofía de la liberación (1977), §2.4}

Una anestesia funciona como una totalidad: \emph{ocupa cada vez más espacio hasta que parece que no hubiera nada más ---ni otro plan, ni otra manera de descansar, ni otra forma de estar con gente---.} \textbf{Dussel dice que eso es falso: siempre se está más allá de ese horizonte.} \emph{Y Tierradentro lo demuestra a escala de pueblo:} \textbf{una comunidad entera cambió su forma de vestir y de comer sin decidirlo ---y después decidió parar---.} \emph{Fíjate en cómo pararon, porque es lo contrario de una prohibición:} \textbf{se dieron cuenta de lo que estaban perdiendo}. \emph{No fue la fuerza de voluntad: fue volver a ver lo que la totalidad les tapaba.}

\textbf{Si esto que uso desapareciera mañana, ¿qué haría en esas horas? ¿Se me ocurre algo?}
\end{softbox}
\linead{\linewidth}\\[1mm]
\linead{\linewidth}

\begin{softbox}{milcMostaza}{milcGris}{Estoicismo · Séneca · Sobre la ira, libro II (c. 45 d.C.) · cuidado interior}
\textit{«El mayor remedio para la ira es la demora.»}

{\footnotesize\itshape\color{milcNegro!70}--- Séneca · Sobre la ira, libro II (c. 45 d.C.)}

Séneca habla de la ira, pero su remedio es exactamente la herramienta que sirve aquí: \textbf{demorar}. \emph{Una anestesia se toma \textbf{rápido} ---esa es la mitad de su poder---:} \textbf{aparece la emoción y en menos de un minuto la mano ya hizo el gesto}, sin que nadie haya decidido nada. \emph{Y lo que rompe eso no es prohibirse ---prohibirse es un pulso que se pierde---:} \textbf{es \emph{meter tiempo en la mitad}}. \emph{Diez minutos entre el impulso y la acción, y en esos diez minutos hacer otra cosa cualquiera.} \textbf{Muchas veces alcanza,} \emph{porque las emociones fluctúan y bajan solas si uno las deja ---que es justo lo que la anestesia no permite comprobar---.}

\textbf{Entre que aparece el impulso y que lo hago, ¿cuánto tiempo pasa?}
\end{softbox}
\linead{\linewidth}\\[1mm]
\linead{\linewidth}

\begin{softbox}{milcTurquesa}{milcGris}{Floridi · lente de la infoesfera}
\textit{«Somos organismos informacionales ---inforgs--- que compartimos un entorno global hecho, en última instancia, de información: la infoesfera.»}

{\footnotesize\itshape\color{milcNegro!70}--- Luciano Floridi · La cuarta revolución (2014)}

\textbf{Hay una anestesia que no cuesta plata, no huele, no deja resaca y está siempre en el bolsillo,} \emph{y por eso es la más difícil de ver: la pantalla}. \textbf{Y tiene una diferencia con todas las demás que conviene entender:} \emph{está \textbf{diseñada} para eso}. \emph{Del otro lado hay gente cuyo trabajo es que no la sueltes, y aprende de lo que miras para mostrarte más de lo mismo.} \textbf{Eso no te vuelve débil ---te vuelve alguien compitiendo contra un sistema que se calibra solo---.} \emph{Y de ahí sale la única estrategia que suele funcionar:} \textbf{no pelear con la fuerza de voluntad sino cambiar el entorno} ---dejar el teléfono en otro cuarto, quitar la aplicación del inicio, cargarlo lejos de la cama---. \emph{La demora de Séneca, puesta en la casa en vez de en la cabeza.}

\textbf{¿Qué cambio en mi cuarto haría que la anestesia me quedara diez minutos más lejos?}
\end{softbox}
\linead{\linewidth}\\[1mm]
\linead{\linewidth}

\begin{infoband}{milcNegro}
\textbf{Nota para el docente.} Las tres son citas textuales y llevan su referencia debajo. Puedes remitir a la obra en clase.
\end{infoband}

\milcestacion{milcVino}{Mi autoevaluación · marca una casilla por fila}

{\small
\setlength{\extrarowheight}{2.2pt}
\begin{tabularx}{\textwidth}{>{\RaggedRight\arraybackslash\bfseries}p{3.0cm}YYY}
\rowcolor{milcVino}\color{white}\bfseries Criterio & \color{white}\bfseries Lo logré & \color{white}\bfseries En proceso & \color{white}\bfseries Necesito apoyo\\[.6mm]
Comprensión del tema & \milcopcion{Explico las ideas con mis palabras.} & \milcopcion{Reconozco las ideas, falta precisión.} & \milcopcion{Necesito repasar lo básico.}\\[.8mm]
\rowcolor{milcGris}Calidad del mapa & \milcopcion{Identifico qué emoción precede a cada anestesia y tengo una alternativa concreta para las tres más frecuentes.} & \milcopcion{Reconozco lo que uso, pero no sé qué estoy evitando con eso.} & \milcopcion{Sigo pensando que si funciona no hay problema.}\\[.8mm]
Aplicación y producto & \milcopcion{Mi producto cumple los criterios.} & \milcopcion{Está armado, con ajustes pendientes.} & \milcopcion{Aún está en construcción.}\\[.8mm]
\rowcolor{milcGris}Reflexión 5 dimensiones & \milcopcion{Respondí cada una con honestidad.} & \milcopcion{Respondí algunas con detalle.} & \milcopcion{Mis respuestas son muy generales.}\\[.8mm]
Triángulo de pensamiento & \milcopcion{Conecté las 3 voces con mi trabajo.} & \milcopcion{Conecté al menos una.} & \milcopcion{No las relaciono con mi proceso.}\\[.8mm]
\end{tabularx}}

\vspace{3mm}
\textbf{Hoy entendí que:}\\[1mm]
\linead{\linewidth}\\[1mm]
\linead{\linewidth}

\textbf{Mi compromiso para la próxima sesión es:}\\[1mm]
\linead{\linewidth}\\[1mm]
\linead{\linewidth}

\begin{softbox}{milcMostaza}{milcGris}{Cita de despedida · Marco Aurelio}
\textit{``El obstáculo en el camino se vuelve el camino. Lo que estorba al camino se vuelve camino.''}\\[1mm]
Cada error de hoy, bien observado, se convierte en pieza del camino que pisarás mañana.
\end{softbox}

\small
\begin{tabularx}{\textwidth}{>{\bfseries}p{3.6cm}Y}
\rowcolor{milcGradeDark!12}Nombre completo: & \linead{10cm}\\
Fecha: & \linead{4cm} \quad Curso/grupo: \linead{4cm}\\
Mi firma: & \linead{9cm}\\
\end{tabularx}
\normalsize

\begin{infoband}{milcGradeDark}
\textbf{Para la próxima sesión.} Dos cosas. \textbf{Primero, prueba \emph{una} sustitución} ---la de la situación más frecuente--- y anota qué pasó: \emph{si la usaste, si funcionó, y si no, por qué no}. \emph{Anota también los diez minutos de demora si alcanzaste a hacerlos.} \textbf{Y si al escribir tu lista te quedaste preocupado ---porque lleva meses, porque aumentó, o porque ya intentaste dejarlo y no pudiste---, habla con un adulto esta semana.} \emph{Usa la ruta del momento 3. Eso no es un fracaso: es lo que corresponde.} \textbf{Segundo, pregunta en tu casa qué llegó y qué se llevó:} averigua con alguien mayor \textbf{si en su pueblo, su vereda o su barrio llegó alguna vez algo que prometía plata rápida} ---un cultivo, un negocio, una bonanza, una mina--- y pregúntale \textbf{qué cambió en la gente} y \textbf{si alguien se dio cuenta a tiempo}. \textbf{Trae:} el resultado de tu sustitución y lo que te contaron. \emph{Vas a oír la historia de Tierradentro contada con otros nombres, y casi siempre con el mismo detalle: nadie recuerda el día en que la cosa se volvió un problema. Recuerdan lo que dejaron de hacer.}
\end{infoband}

\end{document}
